\documentclass[11pt]{article}
\usepackage{icmlko}
\begin{document}

\title{되살아남은 한 방향으로만 간다 --- 용량은 자료를 깨우고 자료는 용량을 못 깨운다}
\author{월드모델 실험실}
\date{노트 520 · 2026년 8월 4일}
\maketitle

\begin{abstract}
노트 519 가 조항을 냈다 --- \textbf{용량이 바뀌면 그 용량에서 닫았던 갈래를
다시 연다}(깊이 $6 \to 12$ 가 나이 감쇠 무게를 $-0.0005$ 에서 $+0.0036$ 로
되살렸다). 그리고 그 조항을 \emph{대칭으로} 읽어 사전 등록했다 --- 이
사이클이 자료 쓰는 법을 네 번 바꿨으니 그 이전에 죽은 \textbf{용량} 후보를
다시 열자고.

\textbf{일곱 중 하나도 안 살아났다.}

\begin{center}
\begin{tabular}{lrrl}
\hline
후보(죽은 노트) & 판 차 & 부호 & $95\%$ 구간 \\
\hline
K $64$(456) & $+0.0005$ & $7/12$ & $[-0.0007, +0.0016]$ \\
깊이 $16$(458) & $-0.0003$ & $5/12$ & $[-0.0011, +0.0006]$ \\
깊이 $24$(458) & $-0.0007$ & $3/12$ & $[-0.0015, +0.0001]$ \\
$lr$ $0.06$(459) & $-0.0021$ & $1/12$ & $\mathbf{[-0.0029, -0.0013]}$ \\
\texttt{it} $400$(473) & $-0.0024$ & $0/12$ & $\mathbf{[-0.0031, -0.0018]}$ \\
$lr$ $0.16$(474) & $-0.0025$ & $1/12$ & $\mathbf{[-0.0037, -0.0013]}$ \\
\texttt{it} $120$(459) & $-0.0032$ & $0/12$ & $\mathbf{[-0.0038, -0.0026]}$ \\
\hline
\end{tabular}
\end{center}

넷은 \textbf{확실히 음수}, 셋은 동점. \textbf{반증 --- 되살아남은 대칭이
아니다.}
\end{abstract}

\begin{figure}[t]
\centering
\includegraphics[width=\textwidth]{figs/f1.pdf}
\caption{\textbf{가} --- 일곱 후보. \textbf{나} --- 한 방향.
\textbf{다} --- 왜 그런가(읽기).}
\end{figure}

\section{조항을 반쪽으로 줄인다}

\begin{quote}
\textbf{조항(정정)} --- \textbf{용량이 바뀌면 그 용량에서 닫았던 \emph{자료
쪽} 갈래를 다시 연다.} 그 역은 이 자료에서 성립하지 않는다 --- 자료 쓰는
법이 바뀌어도 죽은 용량 후보는 죽은 채다.
\end{quote}

\textbf{읽으면} --- 용량은 모형이 \emph{무엇을 표현할 수 있나}를 정한다.
깊이를 $6$ 에서 $12$ 로 올리면 새 상호작용이 표현 가능해지고, 그 위에서
자료 쪽 후보가 \emph{새로} 값을 가질 수 있다. 반대로 자료 쓰는 법은
\emph{어느 행을 보나}만 바꾸고 \textbf{표현할 수 있는 것은 그대로}다 ---
그러니 용량 후보의 값이 안 바뀐다.

이것이 맞다면 예측이 하나 따라온다 --- \textbf{축을 늘리는 변경도(표현할
수 있는 것이 늘어나므로) 자료 쪽 갈래를 되살릴 수 있어야 한다.} 아직 안
쟀다.

\section{덤 --- 용량 넷이 여전히 봉우리다}

이 사이클이 자료 쓰는 법을 네 번 바꿨는데도 깊이 $12$ · \texttt{it} $220$
· $lr$ $0.10$ · $K$ $32$ 가 그대로 최적이다. 셋(깊이 $16 \cdot 24$, $K$
$64$)은 동점이고 넷은 확실히 나쁘다.

\textbf{용량 쪽은 두 해째 닫혀 있다}(노트 $459 \sim 461$ 이 닫았고 이
노트가 다시 확인했다). \textbf{그리고 그 사이 판은 $0.4812 \to 0.4932$ 로
올랐다 --- 전부 자료 쪽에서.}

\section*{이 사이클 결산}

\begin{center}
\begin{tabular}{lr}
\hline
& 판 \\
\hline
사이클 시작 & $0.4812$ \\
\texttt{DOMDROP}(노트 498) & $+0.0045$ \\
\texttt{TRAINCAP}(노트 509) & $+0.0021$ \\
배선 고침(노트 513) & $+0.0023$ \\
\texttt{TRAINW}(노트 514) & $+0.0010$ \\
\texttt{TRAINR}(노트 516) & $+0.0020$ \\
\hline
지금 & $\mathbf{0.4932}$ \\
\hline
\end{tabular}
\end{center}

\textbf{다섯 걸음이 전부 ``어떤 행으로 배울까''이고, 그중 하나는 버그를
고친 것이다.} 노트 $505 \cdot 506$ 이 모형 갈래를 닫은 뒤 이 실험실이 번
$+0.0120$ 에 모형 변경은 \textbf{한 건도 없다.}

\section*{남은 것}

\begin{itemize}
\item \textbf{축을 늘리면 자료 갈래가 되살아나나} --- 위의 예측. 축 추가는
표현력을 늘리므로 용량과 같은 쪽이어야 한다.
\item \textbf{원장에 ``이 채택이 되살릴 수 있는 갈래'' 칸} --- 노트 519 가
제안했다. 사람이 안 잊게 하는 유일한 방법.
\item \textbf{자료 쪽} --- 노트 $506$ 의 ``행 두 배 $= +0.0245$''. 이
사이클이 그 절반을 \emph{있던 행으로} 벌었고 나머지 절반은 \textbf{새 행이
있어야 한다.}
\end{itemize}

\end{document}
